\section*{Capítulo \ref{ch:b1:populations} --- Populações e demografia}

\begin{solution}{exo:b1:populations:1}
\omterm{def:b1:populations:population}{População}: os indivíduos de uma espécie num lugar e num momento,
capazes de cruzar entre si. Densidade: indivíduos por unidade de área
ou de volume. Dispersão: seu padrão espacial --- agregada (uma manada,
uma mancha de urtigas), uniforme (alcatrazes que nidificam a uma
bicada de distância, arbustos de creosoto), aleatória
(dentes-de-leão num gramado).
\end{solution}

\begin{solution}{exo:b1:populations:2}
$l_x$: fração de uma coorte viva à idade $x$. $m_x$: filhas por fêmea
de idade $x$. $R_0 = \sum l_x m_x$: filhas por fêmea recém-nascida ao
longo da vida. $T = \sum x l_x m_x / R_0$: idade média das mães.
$R_0 = 1$: cada fêmea se repõe; a \omterm{def:b1:populations:population}{população} é estável.
\end{solution}

\begin{solution}{exo:b1:populations:3}
Tipo III: cerca de 15 em 1000, \qty{1.5}{\%}. Tipo I: cerca de 700 em
1000, \qty{70}{\%}.
\end{solution}

\begin{solution}{exo:b1:populations:4}
$\dd N/\dd t = rN$ e $\dd N/\dd t = rN(1 - N/K)$. $r$ é a taxa de
crescimento per capita quando os recursos são ilimitados (nascimentos
menos mortes por indivíduo por unidade de tempo); $K$ é a \omterm{thm:b1:populations:logistic}{capacidade
de suporte}, o tamanho em que o crescimento cessa.
\end{solution}

\begin{solution}{exo:b1:populations:5}
$r = \ln(1000)/5 = \qty{1.38}{h^{-1}}$; duplicação $\ln 2/1.38 =
\qty{0.5}{h}$.
\end{solution}

\begin{solution}{exo:b1:populations:6}
$l_x m_x = 0, 0.5, 0.6, 0.2, 0$: $R_0 = 1.3$. $T = (1\times 0.5 +
2\times 0.6 + 3\times 0.2)/1.3 = 2.3/1.3 = 1.77$. $r = \ln 1.3/1.77 =
\qty{0.15}{}$ por unidade de tempo: ela cresce.
\end{solution}

\begin{solution}{exo:b1:populations:7}
$N = 50\times 70/7 = 500$. Só 40 peixes marcados restavam: $N = 40\times
70/7 = 400$; a primeira estimativa era alta demais por um quarto (a
fração de marcados na \omterm{def:b1:populations:population}{população} era menor do que se supôs).
\end{solution}

\begin{solution}{exo:b1:populations:8}
$0.2\times 100\times 0.9 = 18$; $0.2\times 500\times 0.5 = 50$;
$0.2\times 900\times 0.1 = 18$ por ano. Máximo $rK/4 = 50$ por ano, em
$N = 500$.
\end{solution}

\begin{solution}{exo:b1:populations:9}
Dente-de-leão: $r$ (milhares de sementes, anual). Elefante: $K$ (um
filhote em quatro anos, vida longa, cuidado). Bacalhau: $r$ (milhões
de ovos, sem cuidado). Gorila: $K$ (um filhote em quatro anos, décadas
de vida). Mosca-doméstica: $r$ (centenas de ovos, madura em dez dias).
\end{solution}

\begin{solution}{exo:b1:populations:10}
$r = \ln(29/10)/2 = \qty{0.53}{h^{-1}}$; $K \approx 660$. A \qty{8}{h}:
$660/(1 + 65e^{-4.24}) = 660/(1 + 0.94) = 340$, contra 351 contados.
Mais rápido em $K/2 = 330$, atingido em $t = \ln 65/0.53 = \qty{7.9}{h}$,
onde $\dd N/\dd t = rK/4 = 87\times 10^5$ \omterm{def:b1:cell-unit-of-life:cell}{células} por mililitro por
hora.
\end{solution}

\begin{solution}{exo:b1:populations:11}
A regulação exige que a taxa per capita caia quando $N$ sobe, de modo
que $N$ retorne para um equilíbrio depois de uma perturbação; um fator
que mata uma fração fixa altera $N$ mas não a dependência da taxa em
relação a $N$, e por isso, depois dele, a \omterm{def:b1:populations:population}{população} retoma o mesmo
crescimento. O tempo atmosférico limita os insetos por golpear com
frequência bastante, e ao acaso, para que a \omterm{def:b1:populations:population}{população} seja recuada
antes de poder aproximar-se de qualquer teto: uma \omterm{def:b1:populations:population}{população} não
regulada, mantida baixa por perturbações repetidas, cujo tamanho médio
depende da frequência das más estações e não de uma \omterm{thm:b1:populations:logistic}{capacidade de
suporte}.
\end{solution}

\begin{solution}{exo:b1:populations:12}
Duas \omterm{def:b1:populations:population}{populações} do mesmo tamanho, uma com uma base larga de jovens e
outra com classes etárias uniformes, têm futuros diferentes: a
primeira crescerá por décadas mesmo que seu $m_x$ caia agora para o
nível de reposição, porque seus muitos jovens ainda não se
reproduziram (momento demográfico); a segunda não. $N$ registra o
passado; a distribuição etária e o calendário de $m_x$ sobre ela são o
que as equações manipulam, e a pirâmide prevê o tamanho da próxima
geração, ao passo que $N$ nada diz a respeito.
\end{solution}

\begin{solution}{pb:b1:populations:1}
\textbf{1.} $r = \ln(29/10)/2 = \qty{0.53}{h^{-1}}$.
\textbf{2.} $\ln 2/0.53 = \qty{1.3}{h}$.
\textbf{3.} $K \approx 660\times 10^5 = \qty{6.6e7}{}$ \omterm{def:b1:cell-unit-of-life:cell}{células} por
mililitro.
\textbf{4.} $N(t) = 660/(1 + 65e^{-0.53t})$: às 4 h, $660/(1 + 7.8) =
75$ (71 contados); às 8 h, 340 (351); às 12 h, $660/(1 + 0.112) = 594$
(595).
\textbf{5.} Em $K/2 = 330$, atingido em $t = \ln 65/0.53 = \qty{7.9}{h}$;
taxa máxima $rK/4 = 0.53\times 660/4 = 87\times 10^5$ \omterm{def:b1:cell-unit-of-life:cell}{células} por
mililitro por hora.
\textbf{6.} $10e^{0.53\times 12} = 10\times 578 = 5780$: dez vezes os
595 observados.
\textbf{7.} Esgotamento do açúcar (ou do oxigênio); acúmulo de etanol
e de ácidos que envenenam as \omterm{def:b1:cell-unit-of-life:cell}{células}.
\textbf{8.} $q_0 = 0.40$; $q_4 = 1 - 0.34/0.40 = 0.15$: forte
mortalidade inicial e depois uma taxa mais constante --- entre o tipo
II e o III, como na maioria dos grandes mamíferos silvestres.
\textbf{9.} $l_x m_x = 0, 0, 0.312, 0.368, 0.320, 0.272, 0.182, 0.080,
0.018, 0$: $R_0 = 1.55$.
\textbf{10.} $\sum x l_x m_x = 0.624 + 1.104 + 1.280 + 1.360 + 1.092 +
0.560 + 0.144 = 6.16$; $T = 6.16/1.55 = \qty{4.0}{yr}$.
\textbf{11.} $r = \ln 1.55/4.0 = \qty{0.11}{yr^{-1}}$; duplicação em
\qty{6.3}{yr}.
\textbf{12.} $200e^{1.1} = 600$ veados.
\textbf{13.} $N(10) = 600/(1 + 2e^{-1.1}) = 600/1.67 = 360$; passa de
500 quando $(K - N_0)/N_0\,e^{-rt} = 0.2$, isto é, $e^{-rt} = 0.1$, $t = \ln
10/0.11 = 21$ anos.
\textbf{14.} $120\times 150/18 = 1000$ percas.
\textbf{15.} 100 marcadas: $100\times 150/18 = 833$; a primeira
estimativa era alta demais (havia menos peixes marcados soltos do que
se supôs).
\textbf{16.} $rK/4 = 100$ por ano, em $N = 500$.
\textbf{17.} $T_2 = \ln 2/0.4 = \qty{1.7}{yr}$.
\textbf{18.} Captura $hN$ com $h = 0.2$: $rN(1 - N/K) = hN$ dá $N^* =
K(1 - h/r) = 1000\times 0.5 = 500$ e uma captura de \num{100} por ano
--- abaixo da produção máxima sustentável, mas uma fração se ajusta ao
estoque: se o estoque cai à metade, a captura também, e a \omterm{def:b1:populations:population}{população}
volta.
\textbf{19.} Em 500: $0.4\times 500\times 0.5 = 100 < 120$: a
\omterm{def:b1:populations:population}{população} declina. Em 300: $0.4\times 300\times 0.7 = 84 < 120$:
declina mais depressa --- uma vez abaixo de $K/2$, uma colheita fixa
acima da produção máxima leva à extinção.
\textbf{20.} Em $K/2$ o excedente é máximo, mas qualquer erro --- uma
\omterm{def:b1:populations:population}{população} superestimada em um quinto, um ano ruim que reduz $r$ ---
transforma uma retirada sustentável em declínio, e abaixo de $K/2$ o
excedente encolhe junto com a \omterm{def:b1:populations:population}{população}: o equilíbrio é instável ao
excesso de colheita. Colher um pouco acima de $K/2$ deixa uma margem.
\textbf{21.} $R_0$ cai à metade, $0.78 < 1$: a manada encolhe até
ficar abaixo de 400 e o $m_x$ se recuperar --- um fator dependente da
densidade, que regula a manada perto de 400.
\textbf{22.} Não: ele retira \qty{40}{\%} seja qual for o tamanho e
deixa a taxa per capita inalterada; nos anos seguintes a manada volta
a crescer com o mesmo $r$ a partir de um começo mais baixo --- uma
perturbação, não uma regulação.
\textbf{23.} $0.11N(1 - N/600) = 30$: $N^2 - 600N + 163\,600 = 0$,
$N = 300\pm\sqrt{90\,000 - 163\,600}$ --- sem solução real: 30 por ano
excede o excedente máximo $rK/4 = 16.5$, e os lobos levam a manada à
extinção. Com 15 por ano: $N = 300\pm\sqrt{90\,000 -
81\,800} = 300\pm 91$, isto é, 209 e 391; o superior é estável (acima
dele o excedente é menor que 15 e a manada recua; abaixo dele é maior
e ela sobe), o inferior é instável.
\textbf{24.} O número de linces responde ao alimento com atraso: um
aumento de lebres eleva os nascimentos e a \omterm{def:b1:populations:lifetable}{sobrevivência} dos linces ao
longo do ano ou dos dois anos seguintes, e um colapso das lebres só
faz os linces passar fome depois que sua \omterm{def:b1:lipids:triglyceride}{gordura} e os filhotes da
estação se foram.
\textbf{25.} Levedura: $K = \qty{6.6e7}{}$ \omterm{def:b1:cell-unit-of-life:cell}{células} por mililitro, $r =
\qty{0.53}{h^{-1}}$. Veados: $R_0 = 1.55$, $T = \qty{4.0}{yr}$, $r =
\qty{0.11}{yr^{-1}}$.
\end{solution}
